\documentclass[11pt]{amsart}
\usepackage[T1]{fontenc}
\usepackage{lmodern,microtype,amsmath,amssymb,amsthm,mathtools}
\usepackage[margin=1.1in]{geometry}
\usepackage[hidelinks]{hyperref}
\usepackage{xurl,booktabs}

\newtheorem{theorem}{Theorem}[section]
\newtheorem{corollary}[theorem]{Corollary}
\newtheorem{lemma}[theorem]{Lemma}
\newtheorem{proposition}[theorem]{Proposition}
\theoremstyle{remark}
\newtheorem{remark}[theorem]{Remark}
\numberwithin{equation}{section}

\newcommand{\PP}{\mathbb{P}}
\newcommand{\E}{\mathbb{E}}
\newcommand{\Res}{\operatorname{Res}}

\hypersetup{pdftitle={Sharp ratios of SERPT and M-SERPT to the Gittins policy for two-point job sizes},pdfauthor={Yitzchak Shmalo}}

\begin{document}

\title[SERPT and M-SERPT for two-point job sizes]{Sharp ratios of SERPT and M-SERPT to the Gittins policy for two-point job sizes}
\author{Yitzchak Shmalo}
\address{Einstein Institute of Mathematics, The Hebrew University of Jerusalem, Jerusalem, Israel}
\email{yitzchak.shmalo@gmail.com}
\date{}

\subjclass[2020]{60K25, 68M20, 90B22}
\keywords{M/G/1 queue, scheduling, Gittins policy, SERPT, M-SERPT, SOAP, mean response time}

\begin{abstract}
In the M/G/1 queue with unknown job sizes, the Gittins policy minimizes
the mean response time, and the simpler policy M-SERPT (monotonic
shortest expected remaining processing time) is known to be within a
factor $5$ of it. We determine the supremum of this ratio when the
job size takes two values and ties in rank are broken in
first-come first-served (FCFS) order. Over all two-point laws and all
loads below one, the supremum of the ratio of the mean response time of
M-SERPT to that of Gittins is $C_2=1.06017342\ldots$, an algebraic number
defined through the root in $(0,1)$ of $t^5+4t^4-20t^3+12t^2+16t-8$. It
is approached at a load of about $0.686$, as the larger of the two sizes decreases to
the largest value at which FCFS order is optimal, and it is
not attained. For SERPT the corresponding supremum is
$(7-2\sqrt2)/4=1.04289\ldots$. Both results rest on an exact formula for
the excess of the mean response time over that of Gittins.
\end{abstract}

\maketitle

\section{Introduction}

Consider a single-server queue with Poisson arrivals of rate $\lambda$,
independent and identically distributed job sizes with the law of $X$, and
preemption at no cost. The scheduler does not know the size of a job; it
knows only the amount of service the job has received so far, its
\emph{age}. Each policy considered here gives a job a rank that depends
on its age and serves the job of least rank (Section~\ref{sec:model}).
Among such policies the Gittins policy, named after Gittins~\cite{Git89},
minimizes the mean response time; a proof in full generality, with a review
of the earlier proofs, is \cite[Theorem~V.1 and Section~II]{SH21}. Under the policy of shortest expected
remaining processing time (SERPT) the rank of a job is its expected
remaining size~\cite[Definition~2.2]{SHS20}, and Scully, Harchol-Balter and Scheller-Wolf
\cite{SHS20} proposed monotonic SERPT (M-SERPT), whose rank is the
running maximum of the SERPT rank over the ages already reached. They bounded
the ratio of the mean response time of M-SERPT to that of Gittins by a
function of the load $\rho=\lambda\E X<1$ that is at most $5$, for every size
law \cite[Theorem~5.1]{SHS20}. The proof uses the SOAP (schedule ordered
by age-based priority) framework of the same authors \cite{SHS18}, which
gives the mean response time of any policy whose rank is a function of the
age. They also give a family of three-point laws along which the ratio
tends to $2$ \cite[Section~7]{SHS20}, and they conjecture that the
supremum of the ratio over all size laws is $2$
\cite[Conjecture~8.1]{SHS20}. Scully conjectures the factor $2$ both for M-SERPT
and for SERPT \cite[Chapter~11]{Scu22}. Scully also computes that the same
three-point family, with its parameter chosen as a function of the load, drives the
ratio to $2$ for SERPT as well, so that $2$ would be tight for both policies
\cite[Section~11.3]{Scu22}. Whether the ratio for SERPT is bounded by any
constant is left open in~\cite{SHS20}. For the multiserver variant
M-SERPT-$k$ in heavy traffic, Scully, Grosof and Harchol-Balter proved the
factor $2$ under conditions on the tail of the size law
\cite[Theorem~3.2]{SGH21}.

For M-SERPT in the single-server M/G/1 queue the supremum of the ratio over
all size laws therefore lies between $2$ and $5$, and its value remains open. The Gittins rank has no closed form in
general; for two-point laws it has one. The SOAP formula gives the mean
response time of one policy at one size law, and Scully, Harchol-Balter and
Scheller-Wolf evaluate it on the law that takes the sizes $2$ and $14$ with
equal probability \cite[Section~6.4]{SHS18}. The three policies have
different rank functions, so the formula gives a separate expression in the
size law and the load for each, and the supremum of the ratio over all
two-point laws and all loads is that of a quotient of two such expressions
over four parameters. For two-point laws the three policies belong to one
family of policies whose mean response time is a quadratic in one
parameter, and this reduces the problem to a maximization in two
variables. The supremum for M-SERPT is $C_2=1.06017\ldots$ (Theorem~\ref{thm:ms}), so the
ratio $2$ of the three-point family of \cite[Section~7]{SHS20} is not
approached by two-point laws.

Let $\PP(X=a)=p$ and $\PP(X=b)=q=1-p$
with $0<a<b$ and $0<p<1$, and write $d=b-a$ and $\mu=\E X=a+qd$; a job of size $a$ is \emph{short} and a job of
size $b$ is \emph{long}. Write $T(\pi)$ for the response time in steady
state under a policy $\pi$, and
$T_{\mathrm G}$, $T_{\mathrm{MS}}$ and $T_{\mathrm S}$ for the response
time under Gittins, M-SERPT and SERPT; ties in rank are
broken in first-come first-served (FCFS) order. If $pd\le a$, the
three policies serve jobs in FCFS order and the ratio is $1$
(Lemma~\ref{lem:fcfs}). Suppose $pd>a$. For $0\le w<d$ the policy $\pi_w$, defined in
Section~\ref{sec:means} by a rank function, splits every job into stages.
The \emph{initial stage}, the first $a$ units of service, is served in
FCFS order. The last $w$ units of a long job form its \emph{protected
tail}, which once begun runs to completion. The \emph{middle stage}
between them, of length $d-w$, is given to long jobs in FCFS
order, and a later arrival preempts it. The initial stage of a job is
served ahead of all other work except a protected tail already begun. Gittins, M-SERPT and
SERPT schedule jobs exactly as $\pi_w$ does with $w=a/p$, $w=0$ and
$w=\mu$ (Section~\ref{sec:means}), so one computation with the SOAP
formula \cite[Theorem~5.5]{SHS18} compares all three. Their mean response
times are values of one quadratic in $w$, least at $w=a/p$
(Proposition~\ref{prop:regret}):
\begin{equation}\label{eq:regret-intro}
\E T(\pi_w)-\E T_{\mathrm G}=\frac{\lambda pq}{2(1-\lambda a)}
\Bigl(w-\frac ap\Bigr)^2 .
\end{equation}
For M-SERPT the excess $\lambda qa^2/(2p(1-\lambda a))$ does not depend on
$b$, while for fixed $a$, $p$ and $\lambda$ the mean $\E T_{\mathrm G}$
increases with $b$, so the supremum is approached in the limit
$pd\downarrow a$. There the ratio becomes $1+f(p,t)$ with
$t=\lambda a$, where
\begin{equation}\label{eq:f}
f(p,t)=\frac{t(1-p)(p-t)}{(1-t)\bigl(2p-(1-p+p^2)t\bigr)},
\end{equation}
and a resultant of the two critical-point equations reduces the
maximization of $f$ to the quintic
\begin{equation}\label{eq:F}
F(t)=t^5+4t^4-20t^3+12t^2+16t-8 .
\end{equation}
For SERPT, at fixed $p$ and $\rho$, the excess and
$\E T_{\mathrm G}$ grow at the same rate in $b$, and the supremum is the
limit $b/a\to\infty$.

\begin{theorem}\label{thm:ms}
Let ties in rank be broken in FCFS order. The polynomial $F$ of \eqref{eq:F} has exactly one root $t_*$ in
$(0,1)$. With $f$ as in \eqref{eq:f}, let
\[
p_*=\frac{t_*(1-t_*)+\sqrt{t_*^2(1-t_*)^2+t_*(2-t_*^2)}}{2-t_*^2}
\qquad\text{and}\qquad C_2=1+f(p_*,t_*).
\]
Then
\[
\sup\frac{\E T_{\mathrm{MS}}}{\E T_{\mathrm G}}=C_2
=1.06017\,34219\,12767\,96060\,4515\ldots,
\]
the supremum taken over all $0<a<b$, $p\in(0,1)$ and $\lambda>0$ with
$\rho<1$. It is approached with $a=1$, $p=p_*$, $\lambda=t_*$ and
$b=1+(1+\varepsilon)/p_*$ as $\varepsilon\downarrow0$, and it is not
attained.
\end{theorem}

Numerically $t_*=0.45061\,45807\ldots$ and $p_*=0.65713\,91408\ldots$, and
along the family of Theorem~\ref{thm:ms} the load tends to $t_*/p_*=0.68572\,17182\ldots$.
The extremal parameters thus have a moderate load and lie at the
boundary of the region where FCFS is optimal (Lemma~\ref{lem:fcfs} and
Remark~\ref{rem:fcfs}). Along this family the ratio tends to $C_2$ as $pd\downarrow a$ and equals $1$ at
$pd=a$. The ratio is discontinuous at $pd=a$ because of the tie rule. There the M-SERPT
rank of a long job past age $a$ equals $\mu$, the rank of a waiting job,
so FCFS tie breaking lets it finish. For $pd>a$ its rank exceeds
$\mu$, and every later arrival preempts it (compare Remark~\ref{rem:ties}).

\begin{theorem}\label{thm:s}
Let ties in rank be broken in FCFS order. Over the same parameters,
\[
\sup\frac{\E T_{\mathrm S}}{\E T_{\mathrm G}}=\frac{7-2\sqrt2}{4}
=1.04289\,32188\ldots .
\]
The supremum is approached with $p=1/2$, $\rho=2-\sqrt2$ and
$b/a\to\infty$, and it is not attained.
\end{theorem}

FCFS tie breaking is the convention of \cite[Section~2.2]{SHS18} and
\cite[Section~2.1]{SHS20}, and tie breaking changes the answer. With preemptive last-come
first-served tie breaking, the supremum over two-point laws of the ratio of
M-SERPT to Gittins at load $\rho$ is at least $2/(2-\rho)$
(Remark~\ref{rem:ties}), which is $4/3$ at $\rho=1/2$.

Section~\ref{sec:model} recalls the policies and the SOAP formula.
Section~\ref{sec:means} computes the conditional mean response times and
proves \eqref{eq:regret-intro}. Sections~\ref{sec:ms} and~\ref{sec:s}
prove Theorems~\ref{thm:ms} and~\ref{thm:s}. Section~\ref{sec:remarks}
collects remarks on heavy traffic and ties, and
Section~\ref{sec:computations} gives the computer algebra and the
interval arithmetic behind the proofs of Theorems~\ref{thm:ms}
and~\ref{thm:s} and the digits of $C_2$.

\section{Policies and the SOAP formula}\label{sec:model}

A policy of the kind considered here assigns to a job of age $u$ a
\emph{rank} $r(u)$ and always serves the job of least rank, breaking ties
in FCFS order. Throughout, $X$ takes the values $a<b$ with probabilities
$p$ and $q=1-p$; $d=b-a$ and $\mu=\E X=a+qd$ are as in the introduction, and we write
\[
A=1-\lambda a,\qquad
B=\frac{\lambda\E X^2}{2(1-\rho)A}.
\]
Since $\lambda a\le\rho<1$, all of $A$, $1-\rho$ and $B$ are positive.

\emph{SERPT} uses the expected remaining size
$r_{\mathrm S}(u)=\E[X-u\mid X>u]$~\cite[Definition~2.2]{SHS20}, so that $r_{\mathrm S}(u)=\mu-u$ for
$0\le u<a$ and $r_{\mathrm S}(u)=b-u$ for $a\le u<b$. \emph{M-SERPT} uses
$r_{\mathrm{MS}}(u)=\max_{0\le v\le u}r_{\mathrm S}(v)$
\cite[Definition~2.4]{SHS20}. The \emph{Gittins} rank~\cite[Definition~2.1]{SHS20}, the reciprocal of
the Gittins index, is, written with conditional expectations,
\[
r_{\mathrm G}(u)=\inf_{v>u}\frac{\E[\min(X,v)-u\mid X>u]}{\PP(X\le v\mid X>u)} .
\]
For a two-point law only $v=a$ and $v=b$ need to be considered. For
$u<v<a$ no job completes and the quotient is $+\infty$; a
value of $v$ strictly between $a$ and $b$ adds service without adding
completions; a value beyond $b$ gives the same quotient as $v=b$. Hence $r_{\mathrm G}(u)=\min\{(a-u)/p,\ \mu-u\}$ for
$0\le u<a$ and $r_{\mathrm G}(u)=b-u$ for $a\le u<b$.

Since $\E[X-u\mid X>u]$ equals $\mu-u$ for $0\le u<a$ and $b-u$ for
$a\le u<b$, the law of $X$ is new better than used in expectation, that
is, $\E[X-u\mid X>u]\le\mu$ for all $u$, exactly when $d\le\mu$, that is,
when $pd\le a$. For laws with this property Gittins and M-SERPT serve jobs
in FCFS order; this is recalled in~\cite[Section~8.2]{SHS20}, where the case of
Gittins is credited to Aalto, Ayesta and Righter~\cite{AAR09} and the
case of M-SERPT is deduced from their~\cite[Proposition~9]{AAR11}. For two-point laws the statement, with SERPT included, follows directly from the rank functions.

\begin{lemma}\label{lem:fcfs}
If $pd\le a$, then Gittins, M-SERPT and SERPT all serve jobs in FCFS order,
and $\E T_{\mathrm G}=\E T_{\mathrm{MS}}=\E T_{\mathrm S}$.
\end{lemma}

\begin{proof}
The condition $pd\le a$ is equivalent to $d\le\mu$ and to $\mu\le a/p$. So
$r_{\mathrm G}(0)=r_{\mathrm S}(0)=r_{\mathrm{MS}}(0)=\mu$, and for each of
the three policies $r(u)\le\mu$ for all $u$. Indeed, the M-SERPT rank is
the constant $\mu$ on $[0,b)$, since $r_{\mathrm S}\le\mu$ there; the
Gittins and SERPT ranks are nonincreasing from $\mu$ on $[0,a)$ and equal
$b-u\le d\le\mu$ on $[a,b)$. A job that has started therefore never has larger rank
than a job that has not, and waiting jobs all have rank $\mu$. With FCFS
tie breaking, a job once started runs to completion, and jobs start in
FCFS order.
\end{proof}

By the lemma and the optimality of Gittins \cite[Theorem~V.1]{SH21}, FCFS
minimizes the mean response time for these two-point laws among the
policies considered here; by Remark~\ref{rem:fcfs} it does not when
$pd>a$. Righter, Shanthikumar and Yamazaki
proved more: for size laws that are new better than used in expectation,
FCFS minimizes the mean response time among work-conserving policies
that do not use the future \cite{RSY90} (see also \cite[Section~1]{SHS20}).

The SOAP formula \cite[Theorem~5.5]{SHS18} is as follows. Fix a rank function $r$ and a job size $x$, and let
$r^{\mathrm{worst}}(u)=\sup_{u\le v<x}r(v)$ be the worst future rank of a
job of size $x$ at age $u$. For a rank $\vartheta$, the \emph{new work}
$c^{\mathrm{new}}[\vartheta]$ of a job is the service it receives before its
rank first becomes at least $\vartheta$, and $\rho^{\mathrm{new}}[\vartheta]
=\lambda\E c^{\mathrm{new}}[\vartheta]$. This is the work that a later arrival
can do ahead of a job whose worst future rank is $\vartheta$. The \emph{old
work} of a job relative to $\vartheta$ consists of the intervals of age on
which its rank is at most $\vartheta$. The first is the original interval
$X^{\mathrm{old}}_0[\vartheta]$, from age $0$ until the rank first exceeds
$\vartheta$. The others are the recycled intervals
$X^{\mathrm{old}}_i[\vartheta]$, $i\ge1$, on which the rank returns to at
most $\vartheta$. Put
$\rho^{\mathrm{old}}_0[\vartheta]=\lambda\E X^{\mathrm{old}}_0[\vartheta]$. Then,
with $\theta=r^{\mathrm{worst}}(0)$,
\begin{equation}\label{eq:soap}
\E[T\mid X=x]=\frac{\lambda\sum_{i\ge0}\E\bigl(X^{\mathrm{old}}_i[\theta]\bigr)^2}
{2\bigl(1-\rho^{\mathrm{old}}_0[\theta]\bigr)\bigl(1-\rho^{\mathrm{new}}[\theta]\bigr)}
+\int_0^x\frac{du}{1-\rho^{\mathrm{new}}[r^{\mathrm{worst}}(u)]}.
\end{equation}
The conventions \emph{at most} for old work and \emph{at least} for new
work encode FCFS tie breaking \cite[Section~4]{SHS18}.

\section{Conditional mean response times}\label{sec:means}

From now on suppose $pd>a$, and put $h=a/p$. Then
\[
h<\mu<d,
\]
since $\mu-h=q(pd-a)/p$ and $d-\mu=pd-a$. On $[0,a)$ we have
$p(\mu-u)-(a-u)=q(pd-a+u)>0$, so $r_{\mathrm G}(u)=(a-u)/p$ there. The three
rank functions are
\[
\begin{array}{lccc}
 & r_{\mathrm{MS}}(u) & r_{\mathrm G}(u) & r_{\mathrm S}(u)\\ \midrule
0\le u<a & \mu & (a-u)/p & \mu-u\\
a\le u<b & d & b-u & b-u
\end{array}
\]
In each case the rank on $[0,a)$ is at most its value $r(0)$ at age $0$.
On $[a,b)$ the rank is nonincreasing; it exceeds $r(0)$ up to some age $b-w$
and is at most $r(0)$ from then on. Here $r(0)=\mu$ and $w=0$ for M-SERPT,
$r(0)=h$ and $w=h$ for Gittins, and $r(0)=\mu$ and $w=\mu$ for SERPT.
A long job therefore has three stages; a short job has only the first. The initial stage has length $a$ and is
served in FCFS order ahead of every long job in its middle stage. A long job then has a middle
stage of length $d-w>0$, which later arrivals preempt. Last comes a
protected tail of length $w$, which once begun runs to completion.
Two long jobs in their middle stages are served in FCFS order. Under
M-SERPT both have rank $d$ and the tie goes to the earlier arrival. Under
Gittins and SERPT the earlier arrival is at least as old, so its rank $b-u$
is no larger, and at equal ages the tie goes to the earlier arrival. For $0\le w<d$ let $\pi_w$ be the policy with the rank function equal to
$1$ on $[0,a)$, to $2$ on $[a,b-w)$ and to $0$ on $[b-w,b)$. It satisfies
the hypotheses of Lemma~\ref{lem:means} below with $r_0=1$, and it treats
every job in the three stages above, as described in the introduction.

\begin{lemma}\label{lem:means}
Let $0\le w<d$ and $r_0\in\mathbb R$. Let $r$ be a rank function with $r(u)\le r(0)=r_0$
on $[0,a)$, nonincreasing on $[a,b)$, with $r(u)>r_0$ for $a\le u<b-w$ and
$r(u)\le r_0$ for $b-w\le u<b$. Then
\[
\E[T\mid X=a]=a+\frac{\lambda(a^2+qw^2)}{2A},\qquad
\E[T\mid X=b]=B+\frac{b-w}{A}+w .
\]
\end{lemma}

\begin{proof}
Take first $x=a$. Then $\theta=r^{\mathrm{worst}}(0)=r_0$. An earlier job has
rank at most $r_0$ on $[0,a)$, rank above $r_0$ on $[a,b-w)$ and rank at most
$r_0$ on $[b-w,b)$. So $X^{\mathrm{old}}_0[r_0]=a$, there is one recycled
interval, of length $w$ for a long job and $0$ for a short one, and
$\rho^{\mathrm{old}}_0[r_0]=\lambda a$. A later arrival starts at rank
$r_0$, so $c^{\mathrm{new}}[r_0]=0$; likewise $c^{\mathrm{new}}[\vartheta]=0$ for
every $\vartheta\le r_0$, and $r^{\mathrm{worst}}(u)\le r_0$ on $[0,a)$.
Formula \eqref{eq:soap} gives $\lambda(a^2+qw^2)/(2A)+a$.

Take $x=b$. Since $w<d$, the age $a$ lies in $[a,b-w)$, so
$\theta=r^{\mathrm{worst}}(0)=r(a)>r_0$ and every age of an earlier job has
rank at most $r(a)$. Hence $X^{\mathrm{old}}_0[\theta]=X$, there is no
recycled interval, and $\rho^{\mathrm{old}}_0[\theta]=\rho$. A later
arrival has rank at most $r_0$ during its first $a$ units and rank $r(a)$ at
age $a$, so $c^{\mathrm{new}}[\vartheta]=a$ for every $\vartheta\in(r_0,r(a)]$.
For $u<b-w$ the worst future rank $r^{\mathrm{worst}}(u)$ lies in
$(r_0,r(a)]$, and the integrand in \eqref{eq:soap} is $1/A$; for $u\ge b-w$
it is at most $r_0$, and the integrand is $1$. Formula \eqref{eq:soap} gives
$\lambda\E X^2/(2(1-\rho)A)+(b-w)/A+w$.
\end{proof}

For $a=2$, $b=14$, $p=1/2$, Lemma~\ref{lem:means} with $w=\mu=8$ gives
$\E[T\mid X=2]=2+18\lambda/(1-2\lambda)$ and $\E[T\mid X=14]=50\lambda/((1-8\lambda)(1-2\lambda))+6/(1-2\lambda)+8$,
the SERPT formulas of \cite[Section~6.4]{SHS18}.

Write $m(w)=\E T(\pi_w)$ for the mean response time given by Lemma~\ref{lem:means}:
\begin{equation}\label{eq:Tw}
m(w)=p\Bigl(a+\frac{\lambda(a^2+qw^2)}{2A}\Bigr)+q\Bigl(B+\frac{b-w}{A}+w\Bigr).
\end{equation}

\begin{proposition}\label{prop:regret}
Suppose $pd>a$. For $0\le w<d$,
\[
m(w)-\E T_{\mathrm G}=\frac{\lambda pq}{2A}(w-h)^2 .
\]
In particular
\begin{equation}\label{eq:gaps}
\E T_{\mathrm{MS}}-\E T_{\mathrm G}=\frac{\lambda qa^2}{2pA},\qquad
\E T_{\mathrm S}-\E T_{\mathrm G}=\frac{\lambda pq}{2A}\Bigl(\mu-\frac ap\Bigr)^2 .
\end{equation}
\end{proposition}

\begin{proof}
By \eqref{eq:Tw}, $m$ is a quadratic polynomial in $w$ with
$m'(w)=\lambda pqw/A+q(1-1/A)=(\lambda q/A)(pw-a)$ and
$m''=\lambda pq/A$. So $m(w)-m(h)=\lambda pq(w-h)^2/(2A)$, and
$m(h)=\E T_{\mathrm G}$, $m(0)=\E T_{\mathrm{MS}}$, $m(\mu)=\E T_{\mathrm S}$
by Lemma~\ref{lem:means} and the description of the three rank functions
above; all three values of $w$ lie in $[0,d)$.
\end{proof}

The restriction $w<d$ cannot be dropped: at $w=d$ the middle stage
disappears, the policy is FCFS, and \eqref{eq:Tw} no longer applies.

\begin{remark}[FCFS for $pd>a$]\label{rem:fcfs}
FCFS is not optimal when $pd>a$. Write $T_{\mathrm{FCFS}}$ for the response
time under FCFS, the policy with a constant rank function. For it
\eqref{eq:soap} has $X^{\mathrm{old}}_0[\theta]=X$, no recycled interval and
$c^{\mathrm{new}}[\theta]=0$, so
$\E T_{\mathrm{FCFS}}=\mu+\lambda\E X^2/(2(1-\rho))$. Subtracting
$\E T_{\mathrm G}=m(h)$ from \eqref{eq:Tw} gives
\[
\E T_{\mathrm{FCFS}}-\E T_{\mathrm G}=\frac{\lambda q(pd-a)\,Y}{2pA(1-\rho)},
\qquad Y=pb-qa+\lambda a(qb-pd).
\]
Here $Y>0$. If $qb\ge pd$, then $Y\ge pb-qa=(pd-a)+2pa>0$. If $qb<pd$, then
$\lambda a<a/\mu$ because $\rho<1$, so
$\mu Y>\mu(pb-qa)+a(qb-pd)=p(a^2+2qad+qd^2)>0$. Together with
Lemma~\ref{lem:fcfs}, for fixed $a$ and $p$ the value $b=a+a/p$ is the
largest at which FCFS is optimal.
\end{remark}

\section{Proof of Theorem~\ref{thm:ms}}\label{sec:ms}

By Lemma~\ref{lem:fcfs} the ratio is $1$ when $pd\le a$, so we may assume
$pd>a$. By scaling time we may take $a=1$; then $t=\lambda$ satisfies
$0<t<p$, since $\rho=t\mu$ and $\mu>1/p$.

\subsection*{Reduction to the boundary}
Fix $p$ and $t$, and let $d$ vary over the values with $d>1/p$ and
$\rho<1$. By \eqref{eq:gaps}, $\E T_{\mathrm{MS}}-\E T_{\mathrm G}
=tq/(2p(1-t))$ does not depend on $d$. In \eqref{eq:Tw} with $w=h$, the
short-job term does not depend on $d$, while $b-h$, $\E X^2$ and
$1/(1-\rho)$ increase strictly with $d$. So $\E T_{\mathrm G}$ increases
strictly with $d$, and the ratio decreases strictly. Its supremum is the
limit as $d\downarrow1/p$. In that limit
$\mu\to1/p$, $\E X^2\to(1+p-p^2)/p^2$ and $\rho\to t/p$, and by
\eqref{eq:Tw} with $w=h=1/p$ and $b=1+d$,
\[
\E T_{\mathrm G}\to\frac1p+\frac{t(1+p-p^2)}{2p(p-t)},
\]
the FCFS mean response time at $d=1/p$. Dividing,
\begin{equation}\label{eq:boundary}
\sup_{d>1/p}\frac{\E T_{\mathrm{MS}}}{\E T_{\mathrm G}}
=1+\frac{tq/(2p(1-t))}{1/p+t(1+p-p^2)/(2p(p-t))}=1+f(p,t).
\end{equation}
At $d=1/p$ itself the ratio is $1$ by Lemma~\ref{lem:fcfs}, so the
supremum in \eqref{eq:boundary} is a one-sided limit and is not attained.

\subsection*{Maximizing \texorpdfstring{$f$}{f}}
On the triangle $0<t<p<1$ write $f=N/D$ with $N=t(1-p)(p-t)$ and
$D=(1-t)(2p-(1-p+p^2)t)$. Since $1-p+p^2\le1$ we have
$2p-(1-p+p^2)t\ge2p-t\ge p$, so $D>0$ and
\[
0<f(p,t)\le\frac{(1-p)(p-t)}{1-t}\le\min\{1-p,\ p-t\},\qquad
f(p,t)\le\frac{t(1-p)(p-t)}{(1-t)p}\le\frac{t}{1-t}.
\]
Hence $f$ extends continuously by $0$ to the boundary of the closed
triangle $0\le t\le p\le1$, and attains its maximum at an interior
critical point. Direct differentiation gives
\[
N_pD-ND_p=-t(t-1)K,\qquad N_tD-ND_t=(p-1)L,
\]
where
\[
K=(t^2-2)p^2+2t(1-t)p+t,\qquad L=(p^3-2p^2-1)t^2+4pt-2p^2 .
\]
At an interior critical point $K=L=0$. The resultant of $K$ and $L$ with
respect to $p$ is
\begin{equation}\label{eq:res}
\Res_p(K,L)=t^2(t-1)(t^2-t-1)F(t),
\end{equation}
and the first three factors do not vanish for $0<t<1$, so $F(t)=0$. Now
$F(0)=-8$, $F(1)=5$ and
\[
F'(t)=16(1-t)^4+88t(1-t)^3+108t^2(1-t)^2+32t^3(1-t)+t^4>0
\]
on $[0,1]$, so $F$ has exactly one root $t_*$ in $(0,1)$. For fixed
$t\in(0,1)$, $K$ is a quadratic in $p$ with negative leading coefficient
and positive constant term, so it has exactly one positive root, and since
$K|_{p=t}=t(t-1)(t^2-t-1)>0$ and $K|_{p=1}=-(t-1)(t-2)<0$, that root lies in
$(t,1)$. It is the value $p_*$ of the theorem when $t=t_*$. The interior
maximizer is therefore $(p_*,t_*)$, and the supremum of the ratio is
$C_2=1+f(p_*,t_*)$.

Since $t_*<p_*$, the parameters $a=1$, $p=p_*$, $\lambda=t_*$ and
$d=(1+\varepsilon)/p_*$ satisfy $\rho<1$ for small $\varepsilon>0$, and
the ratio tends to $C_2$ as $\varepsilon\downarrow0$ by the reduction
above. Exact rational interval arithmetic
(Section~\ref{sec:computations}) gives
\[
1.060173421912767960604515<C_2<1.060173421912767960604516 .
\]
This proves Theorem~\ref{thm:ms}. \qed

\section{Proof of Theorem~\ref{thm:s}}\label{sec:s}

Again the ratio is $1$ when $pd\le a$, so assume $pd>a$ and take $a=1$.
Write $d=1/p+z$ with $z>0$ and fix $p,\rho\in(0,1)$. Then
$\mu=1/p+qz$, $\lambda=\rho/\mu$ and $\mu-1/p=qz$, so by \eqref{eq:gaps}
\[
\frac{\E T_{\mathrm S}-\E T_{\mathrm G}}{\E T_{\mathrm G}}
=\frac{\rho pq^3(1-\rho)z^2}{\Delta},\qquad
\Delta=2A(1-\rho)\mu\,\E T_{\mathrm G}.
\]
Inserting \eqref{eq:Tw} with $w=1/p$ and $b=1+1/p+z$, and putting
$J=p+q/p$,
\[
\Delta=(1-\rho)\bigl\{(2\mu-\rho)J+2\mu q(1+z)\bigr\}
+\rho q\Bigl\{1+2q\Bigl(\frac1p+z\Bigr)+q\Bigl(\frac1p+z\Bigr)^2\Bigr\}.
\]
As a polynomial in $z$, $\Delta=\Delta_0+\Delta_1z+\Delta_2z^2$ with
\[
\Delta_0=\frac{(1-p\rho)\bigl(2-\rho(1-p+p^2)\bigr)}{p^2},\qquad
\Delta_1=\frac{2q\bigl(2-\rho(1+p^2)\bigr)}{p},\qquad
\Delta_2=q^2(2-\rho),
\]
all positive. Hence
\[
\frac{\E T_{\mathrm S}-\E T_{\mathrm G}}{\E T_{\mathrm G}}
<\frac{\rho pq^3(1-\rho)}{\Delta_2}=pq\,\frac{\rho(1-\rho)}{2-\rho},
\]
and the left side tends to the right side as $z\to\infty$. Now
$pq\le1/4$, with equality at $p=1/2$, and $g(\rho)=\rho(1-\rho)/(2-\rho)$
has $g'(\rho)=(\rho^2-4\rho+2)/(2-\rho)^2$, so on $(0,1)$ it is largest at
$\rho=2-\sqrt2$, where $g=3-2\sqrt2$. The supremum is therefore
$1+(3-2\sqrt2)/4=(7-2\sqrt2)/4$, approached with $p=1/2$,
$\rho=2-\sqrt2$ and $z\to\infty$, and the inequality above is strict for
every finite $b$. \qed

\section{Remarks}\label{sec:remarks}

\begin{remark}[Heavy traffic]
For fixed $a<b$ and $p$ with $pd>a$, the excesses \eqref{eq:gaps} stay bounded
as $\rho\uparrow1$, because $A\to1-a/\mu>0$, while $\E T_{\mathrm G}\ge qB\to\infty$.
So for each such law both ratios tend to $1$ in heavy traffic, and
neither supremum is a heavy-traffic limit: the extremal loads $t_*/p_*$ and $2-\sqrt2$ of Theorems~\ref{thm:ms}
and~\ref{thm:s} lie in $(0,1)$.
\end{remark}

\begin{remark}[Ties]\label{rem:ties}
The M-SERPT rank is constant on intervals of ages, so jobs of different
ages can have equal rank. For
deterministic sizes $X=1$, the M-SERPT rank is constant, so every job has
the same rank. If ties are broken by preemptive last-come first-served,
the mean response time is $1/(1-\rho)$, the mean length of a busy period
started by one job. Gittins has rank $1-u$, strictly decreasing, so it never preempts
the job in service, and with either tie breaking its mean response time is
$1+\rho/(2(1-\rho))$. The ratio is
$2/(2-\rho)$, which equals $4/3$ at $\rho=1/2$, already larger than $C_2$,
and tends to $2$ as $\rho\to1$. Two-point laws come arbitrarily close to
the same ratio $2/(2-\rho)$ under this tie rule. If $pd<a$, the M-SERPT rank is the
constant $\mu$ on $[0,b)$ (proof of Lemma~\ref{lem:fcfs}), so with
preemptive last-come first-served ties its mean response time is
$\mu/(1-\rho)$. Gittins never preempts the job in service and has mean response
time $\mu+\lambda\E X^2/(2(1-\rho))$ with either tie breaking. As $b\downarrow a$ at fixed $\rho$ the
ratio tends to $2/(2-\rho)$.
\end{remark}



\section{Computer-assisted steps}\label{sec:computations}

The identities for the derivatives of $N/D$ (Section~\ref{sec:ms}), the resultant \eqref{eq:res}, the
coefficients of $\Delta$ (Section~\ref{sec:s}) and the display of Remark~\ref{rem:fcfs} with the identity
for $\mu Y$ are identities of rational functions with integer coefficients. They are verified in exact
arithmetic by \texttt{verify.py} and \texttt{identities\_check.py} (Python~3.14, SymPy~1.14 \cite{SymPy}).
The enclosure of $C_2$ uses rational interval arithmetic (the \texttt{fractions} module). For the enclosure of $C_2$, $F$
changes sign between $t_-=0.450614580799471903808270806938$ and $t_+=t_-+2\cdot10^{-30}$, which
therefore contain $t_*$. For every $t\in[t_-,t_+]$, $K$ changes sign in
$p$ between $p_-=0.6571391408919625221961487$ and $p_+=p_-+2\cdot10^{-25}$, which therefore
contain $p_*$. Evaluating $1+f$ in rational interval arithmetic on this
box gives the digits of $C_2$ in Theorem~\ref{thm:ms}.

Along the family of Theorem~\ref{thm:ms} the ratio
$\E T_{\mathrm{MS}}/\E T_{\mathrm G}$ is $1.055082$, $1.060121$ and
$1.060173370$ at $\varepsilon=10^{-1}$, $10^{-3}$ and $10^{-6}$. Along the
family of Theorem~\ref{thm:s} with $d=10$, $10^3$ and $10^6$ the ratio
$\E T_{\mathrm S}/\E T_{\mathrm G}$ is $1.0205$, $1.04259$ and
$1.0428929$.

\section*{Acknowledgments and funding}

The research presented in this paper was supported by the European Research Council (ERC) under the European Union's Horizon Europe research and innovation programme (grant agreement No.~101041711) and by the Simons Foundation as part of the Collaboration on the Mathematical and Scientific Foundations of Deep Learning. It was also supported by Heights Labs, by Convex Nexus Capital and by the Israel Science Foundation (ISF grants 2258/19 and 4101/25).

This work was done by multiple AI agents, Claude (Anthropic) and Codex (OpenAI), connected to Math Brain~\cite{MathBrainRepo} and working together inside the Brain framework~\cite{BrainFramework}. The method is described in~\cite{MathBrainPaper}.

\section*{Declaration of competing interest}

The author holds equity in Heights Labs and in Convex Nexus Capital and has a management role at Convex Nexus Capital.

\section*{Declaration of generative AI and AI-assisted technologies}

This paper is part of a broader research project on proving open problems with AI systems. OpenAI's Codex (GPT-6.1) proposed and proved Theorems~\ref{thm:ms} and~\ref{thm:s} and Proposition~\ref{prop:regret} and wrote \texttt{verify.py} and the simulations, and a separate Codex run (GPT-6.1) audited the proofs with its own computations. Anthropic's Claude (Opus~5.5) re-derived the conditional means from the SOAP formula, the identity of Proposition~\ref{prop:regret}, formula \eqref{eq:boundary} and the bound of Section~\ref{sec:s}. It also wrote \texttt{identities\_check.py}, reran \texttt{verify.py}, and wrote the text. The SOAP formula and the analysis of the sizes $2$ and $14$ are due to Scully, Harchol-Balter and Scheller-Wolf~\cite{SHS18}. M-SERPT, its factor-$5$ bound and the three-point family are due to the same authors~\cite{SHS20}. The factor-$2$ conjecture is theirs~\cite{SHS20} and, for SERPT as well, Scully's~\cite{Scu22}. The Gittins policy is named after Gittins~\cite{Git89}. Its optimality in the M/G/1 queue has several earlier proofs, which Scully and Harchol-Balter review before giving a proof in full generality~\cite[Theorem~V.1 and Section~II]{SH21}. The optimality of FCFS for size laws that are new better than used in expectation is due to Righter, Shanthikumar and Yamazaki~\cite{RSY90}. The coincidence of Gittins with FCFS for such laws is credited in~\cite{SHS20} to Aalto, Ayesta and Righter~\cite{AAR09}, and the same fact for M-SERPT is deduced there from their~\cite[Proposition~9]{AAR11}. All of these results are due to the authors cited. What this paper adds is the identity of Proposition~\ref{prop:regret} and the two sharp constants. The author directed the work and is responsible for the content of the paper and for decisions concerning its dissemination.

\section*{Data and code availability}

The scripts \texttt{verify.py} and \texttt{identities\_check.py}, with
their outputs, are available at
\url{https://github.com/yspennstate/sharp-serpt-gittins-two-point/tree/b19f2dd3cb3413fe47763d2c55de8534509754ed}. The
same repository holds \texttt{simulation\_results.jsonl}, which compares
the exact means with discrete-event simulations at five parameter points.

\begin{thebibliography}{99}

\bibitem{AAR09}
S.~Aalto, U.~Ayesta and R.~Righter, \emph{On the Gittins index in the
M/G/1 queue}, Queueing Syst. \textbf{63} (2009), no.~1--4, 437--458,
doi:10.1007/s11134-009-9141-x.

\bibitem{AAR11}
S.~Aalto, U.~Ayesta and R.~Righter, \emph{Properties of the Gittins index
with application to optimal scheduling}, Probab. Engrg. Inform. Sci.
\textbf{25} (2011), no.~3, 269--288, doi:10.1017/S0269964811000015.

\bibitem{Git89}
J.~C.~Gittins, \emph{Multi-Armed Bandit Allocation Indices}, Wiley,
Chichester, 1989.

\bibitem{RSY90}
R.~Righter, J.~G. Shanthikumar and G.~Yamazaki, \emph{On extremal service
disciplines in single-stage queueing systems}, J. Appl. Probab.
\textbf{27} (1990), no.~2, 409--416, doi:10.2307/3214660.

\bibitem{SGH21}
Z.~Scully, I.~Grosof and M.~Harchol-Balter, \emph{Optimal multiserver
scheduling with unknown job sizes in heavy traffic}, Perform. Eval.
\textbf{145} (2021), 102150, doi:10.1016/j.peva.2020.102150.

\bibitem{SH21}
Z.~Scully and M.~Harchol-Balter, \emph{The Gittins policy in the M/G/1
queue}, arXiv:2111.10703, 2021.

\bibitem{SHS18}
Z.~Scully, M.~Harchol-Balter and A.~Scheller-Wolf, \emph{SOAP: one clean
analysis of all age-based scheduling policies}, Proc. ACM Meas. Anal.
Comput. Syst. \textbf{2} (2018), no.~1, Art.~16, 1--30,
doi:10.1145/3179419.

\bibitem{SHS20}
Z.~Scully, M.~Harchol-Balter and A.~Scheller-Wolf, \emph{Simple
near-optimal scheduling for the M/G/1}, Proc. ACM Meas. Anal. Comput.
Syst. \textbf{4} (2020), no.~1, Art.~11, 1--29, doi:10.1145/3379477.

\bibitem{Scu22}
Z.~Scully, \emph{A new toolbox for scheduling theory}, Ph.D. thesis,
Carnegie Mellon University, technical report CMU-CS-22-132, 2022; revised
June 2025, \url{https://ziv.codes/pdf/scully-thesis.pdf}.

\bibitem{SymPy}
A.~Meurer, C.~P.~Smith, M.~Paprocki, O.~\v{C}ert\'{\i}k,
S.~B.~Kirpichev, M.~Rocklin, A.~Kumar, S.~Ivanov, J.~K.~Moore, S.~Singh,
T.~Rathnayake, S.~Vig, B.~E.~Granger, R.~P.~Muller, F.~Bonazzi,
H.~Gupta, S.~Vats, F.~Johansson, F.~Pedregosa, M.~J.~Curry,
A.~R.~Terrel, \v{S}.~Rou\v{c}ka, A.~Saboo, I.~Fernando, S.~Kulal,
R.~Cimrman and A.~Scopatz,
\emph{SymPy: symbolic computing in Python}, PeerJ Comput. Sci.
\textbf{3} (2017), e103, doi:10.7717/peerj-cs.103.

\bibitem{MathBrainRepo} Y.~Shmalo, \emph{Math Brain}, public repository, 2026, \url{https://github.com/yspennstate/math-brain-public/tree/8c9a6bd478c4ec4082347151b7ae16db50c917dd}.

\bibitem{MathBrainPaper} Y.~Shmalo, \emph{Math Brain: shared, auditable memory for multi-agent mathematical research}, October 2026, doi:10.5281/zenodo.24150837,
\url{https://github.com/yspennstate/math-brain-public/tree/main/papers/math-brain-2026-10-24}.

\bibitem{BrainFramework} Y.~Shmalo, \emph{Brain Framework}, public software repository, 2026, \url{https://github.com/yspennstate/brain-framework-public/tree/5629e644cdd510173ae97e9aa92642dabd22dc57}.
\end{thebibliography}

\end{document}
